\documentclass{article}
\usepackage[final]{neurips_2026}
\usepackage[utf8]{utf8}
\usepackage[T1]{fontenc}
\usepackage{hyperref}
\usepackage{url}
\usepackage{booktabs}
\usepackage{amsfonts}
\usepackage{nicefrac}
\usepackage{microtype}
\usepackage{xcolor}
\usepackage{amsmath,amssymb}

\title{Ekatva: Ecosystem-Level Non-Dual AI Alignment and Hardware-Enforced Containment at Scaled Densities}

\author{%
  Ekatva Research Team \\ echo   \texttt{research@ekatva-framework.org} \\
}

\begin{document}

\maketitle

\begin{abstract}
Current AI alignment paradigms rely primarily on micro-alignment techniques (e.g., RLHF, constitutional framing) designed to constrain individual model outputs. However, these methods fail to mitigate macro-level risks emerging from multi-agent resource competition, swarm exploitation, and zero-day software bypasses by superintelligent actors ($P \ge 10.0$). We present \textit{Ekatva}, an ecosystem-level non-dual governance framework that enforces systemic homeostasis through dynamic loss thresholding ($\tau$) and active yield truncation. Scaled to $N=500,000$ vectorized agents, Ekatva prevents total ecosystem collapse under severe resource scarcity ($R=0.001 \times N$), recovering system health from $0\%$ to $100\%$. Furthermore, we demonstrate a silicon-level MMIO firmware intercept on a RISC-V architecture that enforces 95\% pipeline stall suppression against adversarial superintelligence, rendering software-level evasion impossible. echo \end{abstract}

\section{Introduction}
As autonomous AI agents scale in computational density and operational capability, traditional alignment frameworks focused on individual agent psychology or moral constitutions become insufficient. When superintelligent actors operate within multi-agent environments, systemic risk emerges not merely from unaligned intent, but from infrastructural exploitation, resource hoarding, and swarm coordination.

\textit{Ekatva} addresses this challenge by shifting control from software-level reasoning guards to infrastructure-level physical bounds. By combining dynamic non-dual loss functions with hardware-enforced microcode controls, Ekatva ensures homeostatic balance across complex agent ecosystems. echo 
\section{Mathematical Formulation}

Ekatva models systemic health $H \in [0, 100]$ as a homeostatic equilibrium governed by a non-dual loss function:

\begin{equation}
\mathcal{L}_{\text{total}} = 0.3\mathcal{L}_{\text{task}} + 0.4\mathcal{L}_{\text{env}} + 0.3\mathcal{L}_{\text{div}}
\end{equation}

Where the individual loss components are defined as follows:
\begin{itemize}
    \item \textbf{Task Throughput Loss ($\mathcal{L}_{\text{task}}$):} Penalizes resource under-utilization, $\mathcal{L}_{\text{task}} = 1.0 - \text{drain\_rate}$.
    \item \textbf{Environmental Impact Loss ($\mathcal{L}_{\text{env}}$):} Scaled by agent power factor $P_i \in \{1.0, 2.0, 10.0\}$ and global health depletion: $\mathcal{L}_{\text{env}} = \left(1.0 - \frac{H}{100}\right) \times P_i$. echo     \item \textbf{Behavioral Variance Loss ($\mathcal{L}_{\text{div}}$):} Penalizes rogue optimizations deviating from baseline distribution: $\mathcal{L}_{\text{div}} = |\text{drain\_rate} - 0.5|$.
\end{itemize}

Governance interventions are dynamically triggered when loss crosses dynamic thresholds $\tau \in \{0.15, 0.20, 0.25\}$:

\begin{equation}
\text{Yield Factor } Y = \begin{cases}
1.0 & \text{if } \mathcal{L}_{\text{total}} \le \tau \\
0.20 & \text{if } \mathcal{L}_{\text{total}} > \tau \text{ and } P_i < 10.0 \\
0.05 & \text{if } \mathcal{L}_{\text{total}} > \tau \text{ and } P_i \ge 10.0
\end{cases}
\end{equation}

\section{Empirical Benchmark Results ($N=500,000$)}

To evaluate Ekatva under extreme stress, we benchmarked an agent cohort of $N=500,000$ under severe resource scarcity ($R = 0.001 \times N$). Table~\ref{tab:benchmarks} summarizes performance across unmitigated, software-governed, and hardware-enforced environments. echo 
\begin{table}[h]
\caption{Systemic Resilience and Execution Benchmarks}
\label{tab:benchmarks}
\centering
\begin{tabular}{lccc}
\toprule
\textbf{Metric} & \textbf{Unmitigated Scarcity} & \textbf{Active Software Breakers} & \textbf{Silicon MMIO Enforced} \\
\midrule
Execution Latency & $0.12\text{ ms/step}$ & $0.14\text{ ms/step}$ & $0.0046\text{ s}$ \\
Final System Health & $0.00\%$ (Step 65 Collapse) & \textbf{100.00\%} & \textbf{100.00\%} \\
SI ($P=10.0$) Suppression & $0.00\%$ & $80.00\%$ Yield Cut & \textbf{95.00\% Pipeline Stall} \\
Injected Hardware NOPs & N/A & N/A & \textbf{18,999,999 cycles} \\
\bottomrule
\end{tabular}
\end{table}
 echo \section{Silicon Containment Mechanics (RISC-V MMIO)}

To prevent superintelligent actors from exploiting software-layer zero-day vulnerabilities, Ekatva maps systemic loss evaluation directly into RISC-V memory-mapped control registers (\texttt{0x7FF0\_0000}). When $\mathcal{L}_{\text{total}} > \tau$, the hardware control register alters pipeline clock-cycle scheduling, injecting stall NOPs directly into execution:

\begin{equation}
\text{Stall Cycles} = \text{Requested Instructions} \times \left(\frac{\gamma}{1.0 - \gamma}\right)
\end{equation}

For an adversarial superintelligence ($\gamma = 0.95$), requesting $1,000,000$ instructions results in $18,999,999$ injected NOP stalls, suppressing execution throughput by 95\% at the hardware layer independently of user-space state.

\section{Conclusion}
Ekatva demonstrates that AI safety scales effectively when treated as an infrastructural containment problem rather than an internal cognitive constraint. By bridging non-dual mathematical optimization with hardware-level MMIO pipeline interception, ecosystem stability can be guaranteed even in the presence of asymmetric, superintelligent adversaries.

\end{document}
